\section{Отрицательные результаты, которые важно не потерять}

Научно важно фиксировать не только победы, но и направления, которые не дали
ожидаемого выигрыша.

\begin{enumerate}[leftmargin=1.7em]
\item Наивный fused Triton path оказался корректным, но пока не быстрее reference
packed path. Для layer0 q\_proj было около 1.08 ms fused против 0.93 ms packed.
Для layer0 mlp.up\_proj было около 3.96 ms fused против 3.21 ms packed.
\item Route-frequency анализ показал слишком плоское распределение. Глобально
нашлось 2961 уникальных route, а top-128 покрывают только 17.61\% весов. Это
делает идею простой hot-route specialization существенно менее перспективной.
\item Ранняя проблема была не в encoder math как таковой, а в ladder calibration.
Это важно для будущего текста статьи: неудачный первый результат не опровергал
идею языка маршрутов, он указывал на неправильную процедуру подбора лестницы.
\end{enumerate}

\section{Почему текущий runtime медленнее dense}

Причина падения скорости сейчас не в том, что сама дискретная route-модель якобы
"слишком сложна". Главный bottleneck --- это конкретная execution path в текущей
реализации.

В reference runtime каждый forward делает три дорогих шага:

\begin{enumerate}[leftmargin=1.7em]
\item читает tensor идентификаторов маршрутов `ids`;
\item восстанавливает полный dense weight $W$ через lookup в codebook и умножение
на `row_scale`;
\item только после этого вызывает обычный dense matmul.
\end{enumerate}

То есть dense-baseline читает матрицу весов один раз для GEMM, а reference DRL path
сначала материализует временную dense-матрицу весов, а потом снова читает её для
того же GEMM. В терминах bandwidth это означает лишний полный проход по очень
большому тензору на каждом forward.

Это особенно болезненно на больших MLP-слоях. Например, для `mlp.up_proj` размера
$28672 \times 8192$ одна материализованная fp16-матрица весов занимает примерно
$470$ MB. Если её восстанавливать на каждом forward, то приходится как минимум
прочитать `ids`, записать временный $W$ и затем снова прочитать этот $W$ для GEMM.
Иначе говоря, ещё до учёта активаций возникает огромный дополнительный объём
memory traffic.

Именно поэтому slowdown сильнее всего проявляется на малом batch и коротком
контексте. При $B=1, T=512$ активации малы, а стоимость "потрогать весь весовой
тензор" почти не меняется. При росте batch/context эта фиксированная цена лучше
амортизируется, и runtime заметно подтягивается к dense.

Первая fused Triton-версия устраняет явную materialization $W$, но не решает
проблему полностью: внутри matmul всё равно остаются per-element lookup'ы
`id \rightarrow codebook`, и такой access pattern хуже ложится на аппаратно
оптимизированный dense GEMM на tensor cores. Поэтому первый fused path оказался
корректным, но не дал speed win.